\documentclass[11pt]{article}
\usepackage[margin=25mm]{geometry}
\usepackage{fontspec,unicode-math}
\setmainfont{Latin Modern Roman}
\setmathfont{Latin Modern Math}
\usepackage{amsmath,amssymb,mathtools}
\usepackage{xurl,hyperref}
\hypersetup{colorlinks=true,urlcolor=blue,linkcolor=blue}
\setcounter{secnumdepth}{0}
\setlength{\parindent}{0pt}
\setlength{\parskip}{5pt}
\setlength{\emergencystretch}{3em}
\newcommand{\sref}[2]{\hyperlink{source:#1:#2}{[S#2]}}
\newcommand{\cataloguescope}[1]{\paragraph{Source question.}#1}
\newcommand{\floor}[1]{\left\lfloor#1\right\rfloor}
\begin{document}
\section{Gaussian Mean Testing under Information Constraints}\label{statement}

\subsection{Definitions and mathematical statement}

\paragraph{Shared notation.}Standard number systems are denoted by $\mathbb N,\mathbb Z,\mathbb Q,\mathbb R,\mathbb C$; index ranges, computational models and specialized notation follow the cited problem source and the definitions below.


Gaussian mean testing observes samples from a multivariate normal distribution $N(\mu,\Sigma)$ and distinguishes two specified sets of means. For example, a zero-versus-separated-mean test compares $\mu=0$ with $\lVert\mu\rVert_2\geq\varepsilon$ under a specified covariance; this example explains the terminology without fixing the entire source family. Information constraints restrict which observations or messages reach the decision procedure. Optimality compares achievable sample, communication and error tradeoffs with lower bounds applying to every permitted test.
\paragraph{Statement recorded in the supplied source.}
Determine optimal tests and matching lower bounds for Gaussian mean testing under communication or information constraints. \sref{statement}{2}
\paragraph{Scope qualification.}
The admissible channels and null/alternative classes are part of the cited information-constraint model. The supplied partially solved status is preserved; it is not a claim that every communication-allocation variant remains unresolved.
\paragraph{Source status.}
The supplied catalogue labels this record partially solved. Its source-status metadata was checked against the supplied snapshot on 2026-10-11. This import does not independently establish current open status or mathematical correctness. The catalogue literature assessment is dated 2026-08-22.
\subsection{Short English statement}
Determine the best achievable Gaussian-mean tests when communication or other information restrictions limit what the tester can observe.
\subsection{Sources}
\begin{itemize}
\item[\hypertarget{source:statement:1}{S1}] Ulam AI and the UnsolvedMath Contributors, UnsolvedMath: A Curated Collection of Open Mathematics Problems, catalogue entry. Catalogue curation and metadata: CC BY 4.0. The supplied source question is retained; definitions and the English summary are editorial additions.
\url{https://huggingface.co/datasets/ulamai/UnsolvedMath}.
\item[\hypertarget{source:statement:2}{S2}] New Mathematical Techniques in Information Theory (2023), Oberwolfach Reports, DOI 10.4171/owr/2022/14, pp. 683--707
\url{https://doi.org/10.4171/owr/2022/14}.
\item[\hypertarget{source:statement:3}{S3}] Acharya--Canonne--Tyagi, distributed signal detection
\url{https://proceedings.mlr.press/v125/acharya20b.html}.
\item[\hypertarget{source:statement:4}{S4}] Canonne--Nimitt, distributed Gaussian mean testing
\url{https://arxiv.org/abs/2605.29426}.
\textit{Underlying Oberwolfach Reports material retains its source-specific licence, including CC BY-SA 4.0 where applicable; see THIRD\_PARTY\_NOTICES.md.}
\end{itemize}
\paragraph{Preparation and review.}
Original question attributed to the cited source and UnsolvedMath contributors. Codex (GPT-6; exact serving revision not exposed) assisted with definitions, English summaries, source inspection and formatting. No new human mathematical authorship or independent proof verification is asserted. Catalogue identity checked against the complete supplied registry and existing repository statements; known flagged repetitions and obvious identical conjectures excluded. Full mathematical-equivalence/literature review has not been performed. Statement preparation and final review followed contribution prompts C and D; automated checks verify structure and identity, not open status.
\label{end:statement}
\end{document}
